\subsection{Cancer tumour/normal benchmark}
\label{sec:res-cancer-bench}

We tested whether the phasor spectral-indicator features carry discriminative
power beyond standard baselines, and specifically whether the \emph{phase}
component contributes, on lung adenocarcinoma versus adjacent normal tissue
(GSE10072; 58 tumour, 49 normal; 364 features = top-200 variable probes plus
164 compartment-marker probes). Five feature sets were compared under identical
$\ell_2$-regularised logistic regression with $5\times10$ repeated stratified
cross-validation: the phasor spectral indicators (spectral entropy, Fiedler
gap, participation ratio, mode localisation, Kuramoto $R$); a \emph{phase-free}
ablation of the same panel with $\theta_i\!=\!0$ (so the OCM reduces to the
real amplitude-only form $H^0_{ij}=c_{ij}r_ir_j$); the top-10 principal-component
scores of the $\log_{1+}$ expression matrix; and 10 WGCNA-style module
eigengenes. All feature extraction (phasor phase/amplitude standardisation,
the coupling correlation matrix, PCA, and the WGCNA modules) was fit inside the
training folds only, so the held-out AUROC is leakage-free. Confidence intervals
are 2000-sample bootstraps of the pooled out-of-fold predictions.

Every feature set separates the two classes near the ceiling of the metric
(Fig.~\ref{fig:bench}a): phase-free AUROC $=0.997$ (95\% CI $0.989$--$1.000$), phasor
$=0.998$ ($0.993$--$1.000$), PCA $=0.999$ ($0.996$--$1.000$), and WGCNA eigengenes $=0.998$
($0.990$--$1.000$); accuracies were $0.972$--$0.991$. The decisive test---phasor versus its
own phase-free ablation---shows \textbf{no benefit from the phase}: the mean
per-repeat AUROC difference is $+0.001$ (Wilcoxon signed-rank $p=0.80$, Cohen's
$d_z=0.12$ over 10 repeats). Because the full-sample comparison sits against the
metric ceiling, we repeated the phasor-versus-phase-free contrast at reduced
training sizes ($n_{\mathrm{train}}=16,24,40$; 20 resamples each) to expose any
difference off-ceiling (Fig.~\ref{fig:bench}b). The two remained statistically
indistinguishable at every size (paired Wilcoxon $p=0.26$, $0.27$, $0.055$
respectively; at $n_{\mathrm{train}}=16$ the phase-free variant was in fact
marginally higher). On this
tumour/normal contrast the phasor construction matches strong linear baselines,
but the phase factor $e^{i(\theta_i-\theta_j)}$ adds no measurable
discriminative power.

\begin{table}[t]
\centering
\caption{GSE10072 tumour vs.\ normal: held-out AUROC (pooled out-of-fold, with
bootstrap 95\% CI), accuracy, and $F_1$ under $5\times10$ repeated stratified
CV. The phasor and its phase-free ablation are statistically indistinguishable.}
\label{tab:bench}
\begin{tabular}{lcccc}
\hline
Feature set & AUROC & 95\% CI & Accuracy & $F_1$ \\
\hline
Phase-free (amplitude only) & 0.997 & 0.989--1.000 & 0.972 & 0.974 \\
Phasor (full, with phase)   & 0.998 & 0.993--1.000 & 0.972 & 0.974 \\
PCA (top-10)                & 0.999 & 0.996--1.000 & 0.972 & 0.974 \\
WGCNA eigengenes            & 0.998 & 0.990--1.000 & 0.991 & 0.991 \\
\hline
\end{tabular}
\end{table}

\begin{figure}[t]
\centering
\includegraphics[width=\linewidth]{bench_auroc_comparison.png}
\caption{\textbf{Phasor spectral features match baselines but the phase adds no
discriminative power (GSE10072, 58 tumour / 49 normal).}
(a) Held-out AUROC (pooled out-of-fold point, bootstrap 95\% CI whiskers) under
$5\times10$ repeated stratified CV with leakage-safe feature extraction. All
four feature sets---phase-free amplitude-only indicators, the full phasor
indicators, top-10 PCA, and 10 WGCNA module eigengenes---sit at the ceiling of
the metric; the phasor-vs.\ phase-free increment is not significant.
(b) Learning-curve control at small training sizes, where AUROC is below
ceiling: the phasor (with phase) and phase-free ($\theta=0$) variants remain
statistically indistinguishable at every size (paired Wilcoxon $p\geq0.055$;
error bars are $\pm$SEM over resamples).}
\label{fig:bench}
\end{figure}
